% Partie : logique intuitionniste
% Chapitre : propositions comme types
% Section : des preuves aux termes

\documentclass[../../../include/open-logic-section]{subfiles}

\begin{document}

\olfileid[fr]{int}{pty}{pt}

\olsection{Des \usetoken{p}{derivation} aux termes de preuve}

Passer des preuves en déduction naturelle aux termes de preuve,
c'est-à-dire des preuves de \Log{N2i} aux preuves de \Log{tN2i},
est très simple. Il suffit d'inscrire un terme à gauche
de chaque formule à droite d'un séquent de la dérivation ;
on obtient des séquents $\Gamma \Sequent M : !A$.
Nous disons alors que $M$ \emph{témoigne} de~$!A$
dans le contexte~$\Gamma$.
Les règles de \Log{tN2i} déterminent le terme à inscrire.
Les étiquettes des hypothèses sont fixées, y compris
la variable choisie pour justifier chaque axiome.

\begin{prop}
Dans le fragment propositionnel décrit ci-dessus, toute preuve~$\delta$
de $\Gamma \Sequent !A$ dans \Log{N2i} correspond à une preuve~$\delta'$
de $\Gamma \Sequent N : !A$ dans \Log{tN2i}. Le terme $N$ est déterminé
de manière unique par~$\delta$, une fois ses étiquettes fixées.
\end{prop}

\begin{proof}
On raisonne par récurrence sur la hauteur de~$\delta$.
Si $\delta$ est l'axiome $\Gamma \Sequent !A$, justifié
par l'hypothèse étiquetée $x:!A$ dans le contexte,
alors $\delta'$ est l'axiome $\Gamma \Sequent x:!A$.
Si $\delta$ se termine par une règle d'inférence,
l'hypothèse de récurrence fournit, pour chacune de ses
prémisses, un terme de preuve et une preuve dans \Log{tN2i}
du séquent ainsi étiqueté.
On applique la règle correspondante de \Log{tN2i},
qui détermine le terme de preuve de la conclusion.
\end{proof}

\begin{ex}
Considérons la preuve très simple de
$(!A \land !B) \lif (!B \land !A)$ dans \Log{N1i} :
\[
  \AxiomC{$\Discharge{!A\land !B}{x}$}
  \RightLabel{\Elim\land[2]}
  \UnaryInfC{$!B$}
  \AxiomC{$\Discharge{!A\land !B}{x}$}
  \RightLabel{\Elim\land[1]}
  \UnaryInfC{$!A$}
  \RightLabel{\Intro\land}
  \BinaryInfC{$!B \land !A$}
  \DischargeRule{\Intro\lif}{x}
  \UnaryInfC{$(!A \land !B) \lif (!B \land !A)$}
  \DisplayProof
  \]
Sa présentation dans \Log{N2i} est :
\[
  \Axiom$x : !A\land !B \fCenter !A\land !B$
  \RightLabel{\Elim\land[2]}
  \UnaryInf$x : !A\land !B \fCenter !B$
  \Axiom$x : !A\land !B \fCenter !A\land !B$
  \RightLabel{\Elim\land[1]}
  \UnaryInf$x : !A\land !B \fCenter !A$
  \RightLabel{\Intro\land}
  \BinaryInf$x : !A\land !B \fCenter !B \land !A$
  \RightLabel{\Intro\lif}
  \UnaryInf$\fCenter (!A \land !B) \lif (!B \land !A)$
  \DisplayProof
  \]
Nous attribuons les termes de preuve de haut en bas.
Pour un axiome, le terme est la variable de son contexte,
ici~$x$. Les deux règles $\Elim\land[i]$ produisent
respectivement $\proj{2}{x}$ et~$\proj{1}{x}$.
La règle \Intro{\land} réunit ces termes en
$\pair{\proj{2}{x}}{\proj{1}{x}}$.
Enfin, \Intro{\lif} effectue une $\lambd$-abstraction :
elle lie $x$ et enregistre que cette variable étiquetait
l'hypothèse~$!A \land !B$. On obtient
$\lambd[\typeof{x}{!A \land !B}][\pair{\proj{2}{x}}{\proj{1}{x}}]$.
La preuve dans \Log{tN2i} est donc :
\[
  \Axiom$x : !A\land !B \fCenter x: !A\land !B$
  \RightLabel{\Elim\land[2]}
  \UnaryInf$x : !A\land !B \fCenter \proj{2}{x} : !B$
  \Axiom$x : !A\land !B \fCenter x: !A\land !B$
  \RightLabel{\Elim\land[1]}
  \UnaryInf$x : !A\land !B \fCenter \proj{1}{x} : !A$
  \RightLabel{\Intro\land}
  \BinaryInf$x : !A\land !B \fCenter \pair{\proj{2}{x}}{\proj{1}{x}} : 
    !B \land !A$
  \RightLabel{\Intro\lif}
  \UnaryInf$\fCenter \lambd[\typeof{x}{!A \land
  !B}][\pair{\proj{2}{x}}{\proj{1}{x}}] : (!A \land !B) \lif (!B \land !A)$
  \DisplayProof
  \]
\end{ex}

\end{document}
